\origpage{327--337}

\begin{center}
{\Large\sffamily\bfseries\color{BellaLavenderDeep} EXERCICES}
\end{center}
\vspace{0.8em}

\begin{enumerate}[label=\arabic*.,leftmargin=2.2em,itemsep=1.1em]

\item Inverse de la matrice
\[
A=\begin{pmatrix}
1&2&3&\cdots&n\\
0&1&2&\cdots&n-1\\
\vdots&\ddots&\ddots&\ddots&\vdots\\
0&0&0&\cdots&2\\
0&0&0&\cdots&1
\end{pmatrix}.
\]

\item Calculer le déterminant
\[
\begin{vmatrix}
1^n&2^n&\cdots&(n+1)^n\\
2^n&3^n&\cdots&(n+2)^n\\
\vdots&\vdots&&\vdots\\
(n+1)^n&(n+2)^n&\cdots&(2n+1)^n
\end{vmatrix}.
\]

\item Calculer le déterminant de la matrice $A$ d'ordre $n$, de terme général $a_{ij}$ avec $a_{ii}=a_i+b_i$ et $a_{ij}=b_i$ si $i\ne j$.

\item Soit $A,B$ dans $M_n(\mathbb R)$ et
\[
M=\begin{pmatrix}
A&B&B&\cdots&B&B\\
B&A&B&\cdots&B&B\\
\vdots&\vdots&\ddots&&\vdots&\vdots\\
B&B&B&\cdots&A&B\\
B&B&B&\cdots&B&A
\end{pmatrix}\in M_{np}(\mathbb R).
\]
Soit $E=\{\lambda,\det(A-\lambda B)=0\}$. Trouver une condition nécessaire et suffisante portant sur $E$ pour que $M$ soit inversible et dans ce cas calculer $M^{-1}$.

\item Calculer le déterminant
\[
\begin{vmatrix}
a+b&b+c&a+c\\
a^2+b^2&b^2+c^2&a^2+c^2\\
a^3+b^3&b^3+c^3&a^3+c^3
\end{vmatrix}.
\]

\item Calculer le déterminant dit de Van der Monde :
\[
V(x_1,x_2,\ldots,x_n)=
\begin{vmatrix}
1&1&\cdots&1\\
x_1&x_2&\cdots&x_n\\
x_1^2&x_2^2&\cdots&x_n^2\\
\vdots&\vdots&&\vdots\\
x_1^{n-1}&x_2^{n-1}&\cdots&x_n^{n-1}
\end{vmatrix}.
\]

\item Trouver toutes les matrices $A$ de $M_n(\mathbb R)$ telles que, pour toute matrice $X$ de $M_n(\mathbb R)$ on ait
\[
\det(A+X)=\det A+\det X.
\]

\item Montrer qu'une matrice carrée ayant exactement un élément non nul par ligne et par colonne est inversible.

\item Soit $A_n=(a_{ij})_{1\le i,j\le n}$ avec $a_{ij}=(-1)^{\max(i,j)}$, calculer $\det A_n$.

\item Montrer que $n$ fonctions, $f_1,\ldots,f_n$ de $\mathbb R$ dans $\mathbb C$ sont linéairement indépendantes si, et seulement si, il existe $(x_1,\ldots,x_n)$ dans $\mathbb R^n$ tel que
\[
\det((f_i(x_j)))\ne0.
\]

\item Trouver $(x,y,z)\in\mathbb R^3$ tel que
\[
\frac{2x+y}{x+y-z}=\frac{3x+y-z}{z}=\frac{3z}{2x-y}.
\]

\item Déterminer le polynôme caractéristique de $A\in M_n(\mathbb R)$ avec
\[
A=\begin{pmatrix}
0&0&\cdots&0&1\\
0&0&\cdots&0&2\\
\vdots&\vdots&&\vdots&\vdots\\
0&0&\cdots&0&n-1\\
1&2&\cdots&n-1&n
\end{pmatrix}.
\]

\item Soient $E$ et $F$ deux espaces vectoriels de dimensions respectives $p$ et $q$ sur $\mathbb C$. Soit $f\in L(E)$ et $g\in L(F)$. Calculer le déterminant de
\[
\varphi:L(E,F)\longrightarrow L(E,F)
\]
définie par $\varphi(h)=g\circ h\circ f$.

\item Soient $A$ et $B$ dans $M_n(\mathbb R)$ qui commutent, et telles que
\[
\det(A+B)\ge0.
\]
Montrer que pour tout $p$ de $\mathbb N$, $\det(A^p+B^p)\ge0$.

\item Calculer le déterminant $D_n$ de la matrice carrée d'ordre $n$, $A$, de terme général $a_{ij}$, avec $a_{ii}=2\cos\varphi$, $a_{i,i+1}=a_{i,i-1}=1$ et $a_{ij}=0$ sinon.

\end{enumerate}

\begin{center}
{\Large\sffamily\bfseries\color{BellaLavenderDeep} SOLUTIONS}
\end{center}
\vspace{0.8em}

\begin{enumerate}[label=\arabic*.,leftmargin=2.2em,itemsep=1.2em]

\item Comme $\det A=1$, la matrice est inversible. En détaillant le système $Y=AX$ on constate que
\[
\begin{cases}
y_1=x_1+x_2+\cdots+x_n,\\
y_2=x_2+x_3+\cdots+x_n,\\
\hfill\vdots\\
y_{n-1}-y_n=x_{n-1}+x_n,\\
y_n=x_n,
\end{cases}
\]
d'où l'on tire
\[
\begin{cases}
y_1-y_2-(y_2-y_3)=x_1,\\
y_2-y_3-(y_3-y_4)=x_2,\\
\hfill\vdots\\
y_{n-2}-y_{n-1}-(y_{n-1}-y_n)=x_{n-2},\\
y_{n-1}-2y_n=x_{n-1},\\
y_n=x_n,
\end{cases}
\]
d'où
\[
A^{-1}=\begin{pmatrix}
1&-2&1&0&\cdots&0\\
0&1&-2&1&\cdots&0\\
\vdots&\ddots&\ddots&\ddots&\ddots&\vdots\\
0&0&\cdots&1&-2&1\\
0&0&\cdots&0&1&-2\\
0&0&\cdots&0&0&1
\end{pmatrix}.
\]

\item On note $M$ la matrice générale $m_{ij}=(i+j-1)^n$ pour $1\le i,j\le n+1$.

Pour ne pas faire de jaloux, posons $\alpha_i=i-\dfrac12$, $\beta_j=j-\dfrac12$, on a
\[
m_{ij}=(\alpha_i+\beta_j)^n=\sum_{k=0}^n\binom nk\alpha_i^k\beta_j^{n-k}
=\sum_{r=1}^{n+1}\binom n{r-1}\alpha_i^{r-1}\beta_j^{n+1-r}.
\]
Et avec
\[
\alpha_{i,r}=\binom n{r-1}\alpha_i^{r-1},\qquad
\beta_{r,j}=\beta_j^{n+1-r},
\]
$m_{ij}=\sum_{r=1}^{n+1}\alpha_{i,r}\beta_{r,j}$ devient le terme général de la matrice produit $AB$, avec $A=(\alpha_{i,r})_{1\le i,r\le n+1}$ et $B=(\beta_{r,j})_{1\le r,j\le n+1}$ d'où $\det(M)=\det A\det B$.

Dans la $r$ième colonne de $A$ on factorise $\binom n{r-1}$, donc
\[
\det A=\prod_{r=1}^{n+1}\binom n{r-1}
\begin{vmatrix}
\alpha_1^0&\alpha_1^1&\cdots&\alpha_1^n\\
\alpha_2^0&\alpha_2^1&\cdots&\alpha_2^n\\
\vdots&\vdots&&\vdots\\
\alpha_{n+1}^0&\alpha_{n+1}^1&\cdots&\alpha_{n+1}^n
\end{vmatrix}
=\prod_{r=1}^{n+1}\binom n{r-1}\prod_{1\le i<j\le n+1}(\alpha_j-\alpha_i),
\]
et
\[
\det B=
\begin{vmatrix}
\beta_1^n&\cdots&\beta_{n+1}^n\\
\beta_1^{n-1}&\cdots&\beta_{n+1}^{n-1}\\
\vdots&&\vdots\\
\beta_1^0&\cdots&\beta_{n+1}^0
\end{vmatrix}
=\prod_{1\le i<j\le n+1}(\beta_i-\beta_j).
\]
(Voir l'exercice 6 pour le calcul de ces déterminants de Van der Monde.)

Comme $\alpha_j-\alpha_i=j-i$ et $\beta_i-\beta_j=i-j$ on a finalement
\[
D=\prod_{k=0}^{n}\binom nk\prod_{1\le i<j\le n+1}-(j-i)^2
=(-1)^{\frac{n(n+1)}2}\prod_{k=0}^{n}(n-k)^{k+1}\binom nk.
\]

\item On utilise l'aspect $n$-linéaire alterné du déterminant, fonction des vecteurs lignes. On introduit dans les $n$-uplets $B_i=(b_i,\ldots,b_i)$ et
\[
A_i=(0,\ldots,0,a_i,0,\ldots,0),
\]
$a_i$ en $i$ième position.

On a $\det A=\det(A_1+B_1,\ldots,A_i+B_i,\ldots,A_n+B_n)$ est \emph{a priori} somme de $2^n$ termes, mais les $B_i$ étant tous proportionnels, il ne reste que les $n+1$ termes où figure un $B_i$ au plus.

\[
\det A=\det(A_1,\ldots,A_n)+\sum_{i=1}^n\det(A_1,\ldots,A_{i-1},B_i,A_{i+1},\ldots,A_n)
\]
d'où
\[
\det A=a_1\cdots a_n+\sum_{i=1}^n a_1\cdots a_{i-1}b_i a_{i+1}\cdots a_n.
\]

\item $M$ sera inversible si et seulement si son déterminant est non nul. En soustrayant à chaque colonne (blocs) la dernière, puis en ajoutant toutes les lignes à la dernière on a :
\[
\det M=
\begin{vmatrix}
A-B&0&0&\cdots&0&B\\
0&A-B&0&\cdots&0&B\\
\vdots&\vdots&\ddots&&\vdots&\vdots\\
0&0&0&\cdots&A-B&B\\
B-A&B-A&B-A&\cdots&B-A&A
\end{vmatrix}
\]
\[
=\begin{vmatrix}
A-B&0&0&\cdots&0&B\\
0&A-B&0&\cdots&0&B\\
\vdots&\vdots&\ddots&&\vdots&\vdots\\
0&0&0&\cdots&A-B&B\\
0&0&0&\cdots&0&A+(p-1)B
\end{vmatrix},
\]
d'où
\[
\det M=(\det(A-B))^{p-1}\det(A+(p-1)B),
\]
sera $\ne0$ si et seulement si $\det(A-B)\ne0$ et $\det(A-(1-p)B)\ne0$ d'où
\[
(M\text{ inversible})\Longleftrightarrow(1\text{ et }1-p\text{ ne sont pas dans }E).
\]

On cherche $M^{-1}$ sous la forme
\[
\begin{pmatrix}
A'&B'&\cdots&B'\\
B'&A'&\cdots&B'\\
\vdots&\vdots&\ddots&\vdots\\
B'&B'&\cdots&A'
\end{pmatrix}
\]
avec $A'$ et $B'$ dans $M_n(\mathbb R)$, ce qui conduit aux relations
\[
\begin{cases}
AA'+(p-1)BB'=I_n,\\
BA'+AB'+(p-2)BB'=0,
\end{cases}
\]
\[
\Longleftrightarrow
\begin{cases}
AA'-BA'-AB'+BB'=I_n=(A-B)(A'-B'),\\
AA'+(p-1)(BA'+AB')+[(p-1)(p-2)+p-1]BB'=I_n
\end{cases}
\]
(le $1$er $+(p-1)$ fois le $2$e)
\[
\Longleftrightarrow
\begin{cases}
(A-B)(A'-B')=I_n,\\
(A+(p-1)B)(A'+(p-1)B')=I_n,
\end{cases}
\]
\[
\Longleftrightarrow
\begin{cases}
A'-B'=(A-B)^{-1},\\
A'+(p-1)B'=(A+(p-1)B)^{-1},
\end{cases}
\]
d'où
\[
A'=\frac1p\bigl((p-1)(A-B)^{-1}+(A+(p-1)B)^{-1}\bigr)
\]
et
\[
B'=\frac1p\bigl(-(A-B)^{-1}+(A+(p-1)B)^{-1}\bigr).
\]
\mbox{}\footnote{\textbf{校勘：}原式末项印作“$A+(p-1)B^{-1}$”。由上一行 $A'+(p-1)B'=(A+(p-1)B)^{-1}$ 可知，$-1$ 应作用于整个矩阵 $A+(p-1)B$。}

\item On retrouve dans les colonnes les vecteurs
\[
A=\begin{pmatrix}a\\a^2\\a^3\end{pmatrix},\qquad
B=\begin{pmatrix}b\\b^2\\b^3\end{pmatrix},\qquad
C=\begin{pmatrix}c\\c^2\\c^3\end{pmatrix}
\]
de $K^3$. Le déterminant étant fonction 3 linéaire alternée des vecteurs colonnes, on a
\[
\begin{aligned}
D&=\det(A+B,B+C,C+A)\\
&=\det(A,B,C)+\det(B,C,A)\\
&=2\det(A,B,C)
=2abc\begin{vmatrix}1&1&1\\a&b&c\\a^2&b^2&c^2\end{vmatrix}\\
&=2abc(b-a)(c-a)(c-b),
\end{aligned}
\]
(voir exercice n$^\circ$ 6).

\item C'est un classique qui doit être connu. D'abord, si deux éléments $x_i$ et $x_j$ sont égaux, le déterminant est nul, deux colonnes étant égales.

Si on considère le déterminant développé, c'est un polynôme de degré $n-1$ au plus en chaque variable. Si on considère $x_i$ comme une variable, ce polynôme s'annulera si $x_i=x_j$ pour $j\in\{1,\ldots,n\}\setminus\{i\}$, donc ce polynôme est divisible par chaque $x_i-x_j$ pour $i\ne j$.

Soit alors l'expression
\[
\prod_{1\le i<j\le n}(x_j-x_i).
\]
C'est un polynôme homogène de degré $\dfrac{n(n-1)}2$, (nombre de facteurs), par rapport à l'ensemble des variables, comme $V(x_1,\ldots,x_n)$ qui est une somme de $n!$ termes tous de degré
\[
1+2+3+\cdots+n-1=\frac{n(n-1)}2
\]
(il y a un terme de chaque ligne dans un produit du déterminant). Ces deux polynômes sont proportionnels.

Or le terme qui est : de plus haut degré en $x_n$, et parmi ceux-là, de plus haut degré en $x_{n-1}$, puis en $x_{n-2}$, ... est $x_n^{n-1}x_{n-1}^{n-2}\cdots x_2$ dans les 2 polynômes, (à vérifier) d'où
\[
V(x_1,x_2,\ldots,x_n)=\prod_{1\le i<j\le n}(x_j-x_i).
\]

\item Si $n=1$, la relation devient $a+x=a+x$ pour $a$ et $x$ réels elle est toujours vérifiée.

Soit $n\ge2$. Si $A$ vérifie la propriété, $\det(A+A)=2^n\det A=2\det A$ implique $\det A=0$. Puis, $\forall\lambda\in\mathbb R$,
\[
\det(A-\lambda I_n)=\det A+(-\lambda)^n=(-\lambda)^n :
\]
le polynôme caractéristique de $A$ est scindé, donc $A$ est jordanisable sur $\mathbb R$ (chapitre X, §4). Comme le déterminant est stable par passage à une matrice semblable on remplace $A$ par
\[
A'=\begin{pmatrix}
0&1&0&\cdots&0\\
0&0&1&\ddots&\vdots\\
\vdots&\ddots&\ddots&\ddots&0\\
0&\cdots&0&0&1\\
0&\cdots&\cdots&0&0
\end{pmatrix}
\]
ou une matrice du même type par blocs de Jordan.

Si $A'\ne0$, dans $A'$ il y a des 1 au dessus de la diagonale. Soit $X'$ triangulaire inférieure avec des 1 sur la diagonale et avec $x'_{i,i+1}=1-a'_{i,i+1}$ : on a des 0 au dessus de la diagonale donc $\det A'=\det X'=0$, alors que $A'+X'$ est inversible. La relation $\det(A'+X')=\det A'+\det X'$ est exclue.

Donc seule $A'=0$ convient, d'où seule $A=0$ est solution.

\item Soit $i=\sigma(j)$ l'indice de la ligne tel que, pour $j$ fixé, $a_{ij}\ne0$, (il y a un seul terme de la $j$ième colonne non nul). Comme il n'y a aussi qu'un terme non nul par ligne, $j\ne j'\Rightarrow\sigma(j)\ne\sigma(j')$ donc $\sigma$ est une injection de $\{1,\ldots,n\}$ dans lui-même, c'est-à-dire une bijection.

On a alors
\[
\det A=\varepsilon_\sigma a_{\sigma(1),1}a_{\sigma(2),2}\cdots a_{\sigma(n),n}\ne0
\]
donc $A$ est inversible.

Ou bien $A$ est la matrice de l'isomorphisme envoyant $e_j$ sur $a_{\sigma(j),j}e_{\sigma(j)}$.

\item Si on compare les lignes $L_i$ et $L_{i+1}$ de la matrice, on a
\[
L_i=((-1)^i,\ldots,(-1)^i,(-1)^{i+1},\ldots,(-1)^n),
\]
le changement de signe ayant lieu en $(i+1)$ième position, et de même
\[
L_{i+1}=((-1)^{i+1},\ldots,(-1)^{i+1},(-1)^{i+1},(-1)^{i+2},\ldots,(-1)^n),
\]
d'où l'on tire
\[
L_i+L_{i+1}=(0,\ldots,0,2(-1)^{i+1},2(-1)^{i+2},\ldots,2(-1)^n).
\]

Comme
\[
\det(L_1,L_2,\ldots,L_n)=\det(L_1,L_1+L_2,L_2+L_3,\ldots,L_{n-1}+L_n),
\]
(fonction $n$-linéaire alternée des lignes), et que la nouvelle matrice est triangulaire supérieure, de diagonale $-1,2(-1)^2,2(-1)^3,\ldots,2(-1)^n$ on a
\[
\det A_n=2^{n-1}(-1)^{\frac{n(n+1)}2}.
\]

\item On procède par récurrence sur $n$. Si $n=1$, $(\{f_1\}\text{ libre})\Longleftrightarrow(f_1\ne0)\Longleftrightarrow(\exists x_1,f_1(x_1)\ne0)$.

On suppose le résultat vrai à l'ordre $n-1$. On a :
\[
\begin{aligned}
(\{f_1,\ldots,f_n\}\text{ libre})
&\Longleftrightarrow(\{f_1,\ldots,f_{n-1}\}\text{ libre et }f_n\notin\operatorname{Vect}\{f_1,\ldots,f_{n-1}\})\\
&\Longleftrightarrow\bigl(\exists(x_1,\ldots,x_{n-1})\in\mathbb R^{n-1}\text{ tel que }d_{n-1}\ne0\\
&\qquad\text{ et }f_n\notin\operatorname{Vect}(f_1,\ldots,f_{n-1})\bigr),
\end{aligned}
\]
avec
\[
d_{n-1}=\det[(f_i(x_j))_{1\le i,j\le n-1}].
\]

Soit alors
\[
t\longmapsto D(t)=
\begin{vmatrix}
f_1(x_1)&\cdots&f_1(x_{n-1})&f_1(t)\\
\vdots&&\vdots&\vdots\\
f_{n-1}(x_1)&\cdots&f_{n-1}(x_{n-1})&f_{n-1}(t)\\
f_n(x_1)&\cdots&f_n(x_{n-1})&f_n(t)
\end{vmatrix},
\]
développé par rapport à la dernière colonne, on a
\[
D(t)=d_{n-1}f_n(t)+\sum_{j=1}^{n-1}\alpha_jf_j(t).
\]
Si, $\forall t\in\mathbb R$, $D(t)=0$, comme $d_{n-1}\ne0$, c'est que $f_n\in\operatorname{Vect}\{f_1,\ldots,f_{n-1}\}$, ce qui est exclu, donc $\exists x_n$ tel que $D(x_n)\ne0$. On a bien $\{f_1,\ldots,f_n\}$ libre $\Rightarrow\exists(x_1,\ldots,x_n)\in\mathbb R^n$ tel que $\det((f_i(x_j))_{1\le i,j\le n})\ne0$, pour tout $n$.

La réciproque est évidente, si on a $(x_1,\ldots,x_n)$ dans $\mathbb R^n$ tel que le déterminant des $f_i(x_j)$ soit non nul, et si $\{f_1,\ldots,f_n\}$ était liée, l'une des fonctions serait \BellaCorrection{combinaison linéaire}{combinaison linéaires} des autres d'où une ligne de la matrice combinaison linéaire des autres : c'est exclu.

\item On va résoudre l'exercice en convenant que, si un dénominateur est nul, le numérateur associé l'est aussi. Il sera toujours possible d'exclure les solutions donnant des dénominateurs nuls.

On cherche donc $(x,y,z)$ tel que la matrice
\[
M=\begin{pmatrix}
2x+y&3x+y-z&3z\\
x+y-z&z&2x-y
\end{pmatrix}
\]
soit de rang 1 au plus. Si on note $C_1,C_2,C_3$ les 3 vecteurs colonnes de $M$, la famille $\{C_1,C_2,C_3\}$ est de même rang que $\{C_1+C_2+\frac{C_3}{3},C_2,C_3\}$, (matrice de passage inversible car triangulaire avec des 1 sur la diagonale). Donc $M$ est de même rang que
\[
M'=\begin{pmatrix}
5x+2y&3x+y-z&3z\\
\frac{5x}{3}+\frac{2y}{3}&z&2x-y
\end{pmatrix},
\]
mais aussi de même rang que
\[
M''=\begin{pmatrix}
5x+2y&3x+y-z&3z\\
0&-x-\frac y3+\frac{4z}{3}&2x-y-z
\end{pmatrix},
\]
(si $L_1$ et $L_2$ sont les vecteurs lignes de $M'$, le rang de $\{L_1,L_2\}$ est celui de $\{L_1,L_2-\frac{L_1}{3}\}$).

Vu la forme de $M''$, celle-ci est rang inférieur ou égal à 1 si et seulement si
\[
\begin{cases}
5x+2y=0,\\[0.2em]
\dfrac{3x+y-z}{-x-y/3+4z/3}=\dfrac{3z}{2x-y-z},
\end{cases}
\]
ou
\[
\begin{cases}
5x+2y\ne0,\\
-x-\dfrac y3+\dfrac{4z}{3}=0,\\
2x-y-z=0.
\end{cases}
\]

Le 1er cas équivaut à : $\exists t\in\mathbb R$, $x=2t$, $y=-5t$, et
\[
\frac{t-z}{-t/3+4z/3}=\frac{3z}{9t-z}
\]
\[
\Longleftrightarrow\ \exists t\in\mathbb R,\ x=2t,\ y=-5t,\ z^2+3tz+3t^2=0.
\]
Pour $t,z\in\mathbb R$, cette dernière équation \BellaCorrectionNote{n'a que la solution $t=z=0$}{donne $z=\dfrac{-3\pm\sqrt{21}}2t$}{Le discriminant est $9t^2-12t^2=-3t^2$. Sur $\mathbb R$, il n'y a donc de solution que pour $t=0$, puis $z=0$; la formule imprimée avec $\sqrt{21}$ est algébriquement erronée.}, qui donne des dénominateurs nuls et doit être exclue.

Le 2e cas équivaut à :
\[
5x+2y\ne0,\qquad
\begin{cases}3x+y=4z,\\2x-y=z,
\end{cases}
\]
\[
\Longleftrightarrow\ 5x+2y\ne0,\quad3x+y=4z,\quad5x=5z
\Longleftrightarrow x=z=y\ne0.
\]

\item On développe $P_n(X)=\det(A-XI_n)$ suivant la dernière ligne en :
\[
\begin{aligned}
P_n(X)
&=\sum_{i=1}^{n-1}(-1)^{n+i}i\,(-X)^{i-1}(-1)^{n-i+1}i\,(-X)^{n-i-1}\\
&\quad +(n-X)(-X)^{n-1}\\
&=\sum_{i=1}^{n-1}(-1)^{n-1}i^2X^{n-2}+(-1)^{n-1}X^{n-1}(n-X)\\
&=(-1)^nX^{n-2}\left(X^2-nX-\sum_{i=1}^{n-1}i^2\right).
\end{aligned}
\]

\item Soit $\mathcal B=\{e_1,\ldots,e_p\}$ une base de $E$ et $\mathcal C=\{\varepsilon_1,\ldots,\varepsilon_q\}$ une base de $F$. On note $A$ et $B$ les matrices, de termes généraux $a_{ij}$ et $b_{ij}$, de $f$ et $g$ respectivement dans ces bases.

Soit $h_{ij}$ l'application linéaire de $E$ dans $F$ définie par : $h_{ij}(e_i)=\varepsilon_j$ et pour tout $k\ne i$, $h_{ij}(e_k)=0$.

On a
\[
\begin{aligned}
g\circ h_{ij}\circ f(e_k)
&=g\circ h_{ij}\left(\sum_{u=1}^p a_{uk}e_u\right)
=g(a_{ik}\varepsilon_j)\\
&=a_{ik}\sum_{v=1}^q b_{vj}\varepsilon_v.
\end{aligned}
\]
Comme $\varepsilon_v=h_{kv}(e_k)$, on a
\[
(g\circ h_{ij}\circ f)(e_k)=\left(\sum_{v=1}^q a_{ik}b_{vj}h_{kv}\right)(e_k).
\]
Or, pour tout $l\ne k$, $h_{lv}(e_k)=0$, donc le second membre est encore
\[
=\left(\sum_{l=1}^p\sum_{v=1}^q a_{il}b_{vj}h_{lv}\right)(e_k).
\]
Cette égalité étant valable pour tout $e_k$ de $\mathcal B$, on a
\[
\varphi(h_{ij})=\sum_{l=1}^p\sum_{v=1}^q a_{il}b_{vj}h_{lv}.
\]

Si on indexe la base canonique de $L(E,F)$ dans l'ordre :
\[
h_{11},h_{12},\ldots,h_{1q};\ h_{21},h_{22},\ldots,h_{2q};\ \ldots;\ h_{p1},h_{p2},\ldots,h_{pq},
\]
la matrice de $\varphi$ est la matrice bloc $C$ avec
\[
C=\begin{pmatrix}
a_{11}B&a_{21}B&\cdots&a_{p1}B\\
a_{12}B&a_{22}B&\cdots&a_{p2}B\\
\vdots&\vdots&&\vdots\\
a_{1p}B&a_{2p}B&\cdots&a_{pp}B
\end{pmatrix}.
\]

En prenant une base $\mathcal B$ de $E$ telle que $A$ soit triangulaire supérieure, la matrice $C$ devient triangulaire par blocs et
\[
\det C=\prod_{i=1}^p\det(a_{ii}B).
\]
Comme $B$ est carrée d'ordre $q$, $\det(a_{ii}B)=(a_{ii})^q\det B$ et on obtient
\[
\det C=\left(\prod_{i=1}^p a_{ii}\right)^q(\det B)^p
=(\det A)^{\dim F}(\det B)^{\dim E}.
\]
Finalement
\[
\det\varphi=(\det f)^{\dim F}(\det g)^{\dim E}.
\]

\item Si $p=2q$ est pair, comme $A$ et $B$ commutent,
\[
A^{2q}+B^{2q}=(A^q+iB^q)(A^q-iB^q).
\]
Comme $A$ et $B$ sont à coefficients réels, les matrices $A^q+iB^q$ et $A^q-iB^q$ sont complexes conjuguées donc
\[
\det(A^{2q}+B^{2q})=|\det(A^q+iB^q)|^2
\]
est réel positif.

Si $p=2q+1$ est impair, les $2q+1$ racines de $x^{2q+1}+1=0$ s'écrivent
\[
-1,\alpha_1,\ldots,\alpha_q,\overline\alpha_1,\ldots,\overline\alpha_q,
\]
et, comme $A$ et $B$ commutent, on a
\[
A^{2q+1}+B^{2q+1}=(A+B)\prod_{j=1}^q(A+\alpha_jB)(A+\overline\alpha_jB),
\]
donc
\[
\det(A^{2q+1}+B^{2q+1})=\det(A+B)\prod_{j=1}^q|\det(A+\alpha_jB)|^2
\]
est encore positif puisque $\det(A+B)$ l'est.

\item En développant $D_n$ par rapport à la première ligne on obtient
\[
D_n=(2\cos\varphi)D_{n-1}-
\begin{vmatrix}
1&1&0&\cdots&0\\
0&2\cos\varphi&1&\ddots&\vdots\\
0&1&2\cos\varphi&\ddots&0\\
\vdots&\ddots&\ddots&\ddots&1\\
0&\cdots&0&1&2\cos\varphi
\end{vmatrix},
\]
le déterminant étant d'ordre $n-1$, on le développe par rapport à la 1re colonne, et pour $n\ge3$ on obtient
\[
D_n=2\cos\varphi\,D_{n-1}-D_{n-2},
\]
avec $D_1=2\cos\varphi$ et $D_2=4\cos^2\varphi-1$. On a une suite récurrente linéaire d'ordre 2, d'équation caractéristique (voir le chapitre sur les suites),
\[
r^2-2r\cos\varphi+1=0,
\]
de zéros $r=e^{i\varphi}$ et $e^{-i\varphi}$. Pour $\varphi\ne0\ (\pi)$ les zéros sont distincts donc il existe deux coefficients $\lambda$ et $\mu$ tels que
\[
D_n=\lambda e^{in\varphi}+\mu e^{-in\varphi}.
\]
On a en particulier, ($n=1$ et 2)
\[
D_1=\lambda e^{i\varphi}+\mu e^{-i\varphi}=2\cos\varphi,
\]
et
\[
D_2=\lambda e^{2i\varphi}+\mu e^{-2i\varphi}=4\cos^2\varphi-1.
\]
On en déduit
\[
\begin{cases}
\mu(1-e^{-2i\varphi})=2\cos\varphi\,e^{i\varphi}-4\cos^2\varphi+1,\\
\lambda(1-e^{2i\varphi})=2\cos\varphi\,e^{-i\varphi}-4\cos^2\varphi+1,
\end{cases}
\]
soit encore
\[
\begin{cases}
\mu e^{-i\varphi}2i\sin\varphi=1-2\cos^2\varphi+2i\cos\varphi\sin\varphi=-\cos2\varphi+i\sin2\varphi,\\
\lambda e^{i\varphi}(-2i\sin\varphi)=1-2\cos^2\varphi-2i\sin\varphi\cos\varphi=-\cos2\varphi-i\sin2\varphi,
\end{cases}
\]
soit
\[
2i\mu e^{-i\varphi}\sin\varphi=-e^{-2i\varphi}
\qquad\text{et}\qquad
-2i\lambda e^{i\varphi}\sin\varphi=-e^{2i\varphi},
\]
d'où
\[
\mu=\frac{e^{i(\frac\pi2-\varphi)}}{2\sin\varphi}
\qquad\text{et}\qquad
\lambda=\frac{e^{i(\varphi-\frac\pi2)}}{2\sin\varphi}
\]
donc
\[
D_n=\frac1{2\sin\varphi}\left(e^{i((n+1)\varphi-\frac\pi2)}+e^{-i((n+1)\varphi-\frac\pi2)}\right)
=\frac{\cos((n+1)\varphi-\frac\pi2)}{\sin\varphi},
\]
soit encore
\[
D_n=\frac{\sin(n+1)\varphi}{\sin\varphi}.
\]
Il en résulte, par continuité de $D_n$ en $\varphi$, (expression polynômiale en $\cos\varphi$), que si $\varphi$ tend vers 0, $D_n$ tend vers $(n+1)$ et si $\varphi$ tend vers $\pi$, (modulo $2\pi$), avec $\varphi=\pi+t$,
\[
\frac{\sin(n+1)\varphi}{\sin\varphi}
=\frac{(-1)^{n+1}\sin(n+1)t}{-\sin t}\underset{t\to0}{\longrightarrow}(-1)^n(n+1)
\]
et finalement, pour $\varphi\ne0\ (\pi)$,
\[
D_n=\frac{\sin(n+1)\varphi}{\sin\varphi},
\]
pour $\varphi\equiv0\ (2\pi)$, $D_n=n+1$ et pour $\varphi\equiv\pi\ (2\pi)$,
\[
D_n=(-1)^n(n+1).
\]
\mbox{}\footnote{\textbf{校勘：}原式末行印作“$D_n(-1)^n(n+1)$”，缺少等号。}

\end{enumerate}
